\chapter[Conclusão]{Conclusão}

Esta versão inicial consolidou as quatro entregas anteriores na estrutura do modelo de monografia da Faculdade de Computação. O conteúdo da E1 foi reorganizado na Introdução, com separação explícita entre motivação, problema, hipótese, objetivos e contribuições esperadas. Essa mudança eliminou a organização por parágrafos do enunciado e tornou mais clara a relação entre a relevância do tema, a questão investigada e os critérios necessários para verificar o objetivo.

A fundamentação da E2 foi revisada para apresentar os conceitos na ordem em que são utilizados pela pipeline: análise seminal assistida por computador, detecção, rastreamento, fluxo óptico, predição, dados e métricas. Os textos de caráter didático solicitados exclusivamente nas entregas anteriores foram removidos. Os trabalhos analisados na E3 passaram a compor um capítulo próprio, sem os resumos de capítulos do livro e sem informações conjunturais que não contribuem diretamente para a comparação científica. A tabela comparativa foi mantida e atualizada para evidenciar a lacuna relacionada à integração entre rastreamento, movimento aparente e predição individual.

O método proveniente da E4 foi ajustado para manter correspondência direta com a hipótese. Foram detalhadas as divisões dos dados, as linhas de base, a ablação com e sem fluxo, os cenários com trajetórias anotadas e estimadas e os critérios de avaliação. Essa revisão busca reduzir ambiguidades, evitar vazamento de informação e distinguir o movimento aparente medido na imagem de uma velocidade física do fluido. A declaração sobre o uso de Inteligência Artificial foi mantida somente no capítulo de Método, conforme solicitado para esta entrega.

\section{Estado do desenvolvimento e continuidade}

Na revisão de 8 de setembro de 2026, o projeto permanecia alinhado à hipótese e não exigia reinício dos experimentos. A organização das fontes, a definição das divisões e a busca de threshold v3 no treino já haviam avançado além do planejamento inicial. A triagem e o refinamento estavam concluídos e conferidos, com T219/o0/c2 e T218/o0/c2 como finalistas de treino. A auditoria geral registrou 3.248 arquivos conferidos, 17.753 verificações e aprovação em 502 testes automatizados, sem incluir as verificações opcionais de aprendizagem. Esses resultados atestam o alcance das verificações executadas, não a validação científica de todos os módulos nem a hipótese sobre fluxo óptico.

Naquela revisão, o próximo marco previsto era comparar as duas finalistas nos quatro vídeos completos de validação, conforme protocolo prospectivo, totalizando 11.700 avaliações de configuração--quadro. Os controles para quantidade exata de quadros, rejeição de anotações malformadas e integridade das entradas ainda precisavam ser concluídos e registrados. Esse marco é histórico: a bateria e sua conferência foram realizadas posteriormente, como descrito a seguir. Os resultados anteriores permanecem preservados e separados dos contratos novos.

Após essa preparação, a validação completa selecionou T218/o0/c2, com F1 macro de 0,658959, contra 0,657819 de T219/o0/c2. As 11.700 avaliações e a conferência independente foram concluídas; a diferença foi pequena e cada configuração teve F1 maior em dois vídeos. Na conclusão específica daquela bateria ainda não havia congelamento. Posteriormente, T218/o0/c2 foi fixada somente para desenvolvimento; isso não liberou teste, folds ou aplicação, nem demonstrou contribuição do fluxo à predição.

\subsection{Atualização experimental de 11 de setembro de 2026}

A busca clássica entre seis famílias foi concluída e conferida, com 43 configurações nos mesmos 576 quadros dos 12 vídeos de treino, totalizando 24.768 avaliações. Blob obteve o maior F1 macro da grade, 0,791005, enquanto a referência T218 obteve 0,775634. Blob teve F1 maior em sete vídeos e T218 em cinco; a ordem média não é uniforme entre vídeos. A grade limitada, o esforço desigual de busca e o ajuste anterior de T218 impedem interpretar essa triagem como superioridade geral. As 11 finalistas retidas ainda não constituem promoção ou escolha final da pipeline. O protocolo, a revisão operacional e o alcance da conferência estão descritos na Seção~\ref{sec:busca-classicos-20260911}.

A preparação das trajetórias individuais e os baselines de predição sobre GT já estão concluídos. O benchmark compacto de fluxo também foi conferido, mas a hipótese permanece aberta: não há ainda ADE/FDE com fluxo ou HOTA real da comparação entre rastreadores. Após o refinamento local e a preparação do dataset descritos a seguir, a próxima sequência compreende validação completa das alternativas clássicas, treinamento e comparação de YOLO, conclusão do contrato de rastreamento e avaliação dos métodos nas mesmas detecções. Depois de comparar e validar as combinações pertinentes, a predição deverá receber trajetórias efetivamente estimadas, preservando trocas de identidade, perdas e cobertura. A ablação controlada sobre GT pode avançar em paralelo; não substitui a avaliação do sistema automático. Todas as entradas continuarão limitadas ao instante da previsão, e ADE incluirá todos os passos de cada horizonte.

A hipótese ainda não foi testada por esses resultados de detecção. A futura análise deverá declarar a exposição histórica do teste e a limitação de apenas quatro vídeos no holdout, sem prometer significância de 5\% pelo Wilcoxon bilateral exato. Diferenças por vídeo e cobertura das janelas serão essenciais à interpretação. A comparação controlada poderá sustentar, limitar ou não sustentar a contribuição do fluxo; qualquer dessas conclusões dependerá da evidência efetivamente obtida.


\subsection{Avanço após o refinamento local}

O refinamento das cinco famílias clássicas foi concluído e conferido em 45 configurações únicas, com 25.920 avaliações nos mesmos quadros de treino. As dez finalistas retidas ainda não constituem uma escolha final. A paridade dos dez pais, os hashes e a reconstrução das métricas apoiam a consistência dos resultados, mas não substituem a validação em vídeos completos ou a avaliação de rastreamento e predição. O dataset YOLO de 23.316 pares de treino e validação também foi materializado e conferido; não houve treinamento nesta etapa. As regras, o alcance e as limitações desses avanços estão na Seção~\ref{sec:refinamento-classicos-20260911}.

Espera-se que os resultados permitam identificar não apenas a melhor configuração média, mas também as situações em que cada módulo falha. A versão final deverá substituir esta descrição de expectativas por resultados experimentais, discussão das limitações, verificação dos objetivos e síntese das contribuições efetivamente obtidas.
